\import{ordinal.tex}
\import{set/inductive-definition.tex}

\section{Proto-surreal sign sequences}

A proto-surreal of length $\alpha$ is represented at the logical level by a
predicate on sets.  Only its values below $\alpha$ matter.  This mirrors
Megalodon's Part 2 representation and deliberately keeps logical predicates
distinct from the set codes introduced later.

\begin{definition}\label{pno_eq}
    $\pnoEqOp :
        \Set\to(\Set\to\Prop)\to(\Set\to\Prop)\to\Prop
        =\fun{\alpha:\Set}{
            \fun{p,q:\Set\to\Prop}{
                \forall \beta:\Set.\ \beta\in\alpha\implies
                    (\Apply{p}{\beta}\iff\Apply{q}{\beta})}}$.
\end{definition}

\begin{proposition}\label{pno_eq_reflexive}
    $\forall \alpha:\Set.\ \forall p:\Set\to\Prop.
        \Apply{\pnoEqOp}{\alpha,p,p}$.
\end{proposition}
\begin{proof}
    Fix $\alpha,p$.
    Follows by \obv{pno_eq}.
\end{proof}

\begin{proposition}\label{pno_eq_symmetric}
    $\forall \alpha:\Set.\ \forall p,q:\Set\to\Prop.
        \Apply{\pnoEqOp}{\alpha,p,q}\implies
        \Apply{\pnoEqOp}{\alpha,q,p}$.
\end{proposition}
\begin{proof}
    Fix $\alpha,p,q$.
    Assume $\Apply{\pnoEqOp}{\alpha,p,q}$.
    Follows by \obv{pno_eq}.
\end{proof}

\begin{proposition}\label{pno_eq_transitive}
    $\forall \alpha:\Set.\ \forall p,q,r:\Set\to\Prop.
        (\Apply{\pnoEqOp}{\alpha,p,q}\land
         \Apply{\pnoEqOp}{\alpha,q,r})
            \implies
         \Apply{\pnoEqOp}{\alpha,p,r}$.
\end{proposition}
\begin{proof}
    Fix $\alpha,p,q,r$.
    Assume $(\Apply{\pnoEqOp}{\alpha,p,q}\land
        \Apply{\pnoEqOp}{\alpha,q,r})$.
    Follows by \obv{pno_eq}.
\end{proof}

\begin{proposition}\label{pno_eq_restrict}
    Let $\alpha$ be an ordinal.
    Let $\beta\in\alpha$.
    Then $\forall p,q:\Set\to\Prop.
        \Apply{\pnoEqOp}{\alpha,p,q}\implies
        \Apply{\pnoEqOp}{\beta,p,q}$.
\end{proposition}
\begin{proof}
    Fix $p,q$.
    Assume $\Apply{\pnoEqOp}{\alpha,p,q}$.
    $\beta$ is an ordinal by \obv{elem_of_ordinal_is_ordinal}.
    $\beta\subseteq\alpha$ by \obv{ordinal_elem_implies_subseteq}.
    Follows by \obv{pno_eq,subseteq}.
\end{proof}

\begin{definition}\label{pno_lt_initial}
    $\pnoLtInitialOp :
        \Set\to(\Set\to\Prop)\to(\Set\to\Prop)\to\Prop
        =\fun{\alpha:\Set}{
            \fun{p,q:\Set\to\Prop}{
                \exists \beta:\Set.\ \beta\in\alpha\land
                    \Apply{\pnoEqOp}{\beta,p,q}\land
                    \lnot\Apply{p}{\beta}\land\Apply{q}{\beta}}}$.
\end{definition}

\begin{proposition}\label{pno_lt_initial_intro}
    Let $\beta\in\alpha$.
    Then $\forall p,q:\Set\to\Prop.
        (\Apply{\pnoEqOp}{\beta,p,q}\land
         \lnot\Apply{p}{\beta}\land\Apply{q}{\beta})
        \implies
        \Apply{\pnoLtInitialOp}{\alpha,p,q}$.
\end{proposition}
\begin{proof}
    Fix $p,q$.
    Assume
        $(\Apply{\pnoEqOp}{\beta,p,q}\land
          \lnot\Apply{p}{\beta}\land\Apply{q}{\beta})$.
    Follows by \obv{pno_lt_initial}.
\end{proof}

\begin{proposition}\label{pno_lt_initial_irreflexive}
    $\forall \alpha:\Set.\ \forall p:\Set\to\Prop.
        \lnot\Apply{\pnoLtInitialOp}{\alpha,p,p}$.
\end{proposition}
\begin{proof}
    Fix $\alpha,p$.
    Suppose $\Apply{\pnoLtInitialOp}{\alpha,p,p}$.
    Follows by \auto{pno_lt_initial}.
\end{proof}

\begin{proposition}\label{pno_lt_initial_extend}
    Let $\alpha$ be an ordinal.
    Let $\beta\in\alpha$.
    Then $\forall p,q:\Set\to\Prop.
        \Apply{\pnoLtInitialOp}{\beta,p,q}\implies
        \Apply{\pnoLtInitialOp}{\alpha,p,q}$.
\end{proposition}
\begin{proof}
    Fix $p,q$.
    Assume $\Apply{\pnoLtInitialOp}{\beta,p,q}$.
    $\beta$ is an ordinal by \obv{elem_of_ordinal_is_ordinal}.
    $\beta\subseteq\alpha$ by \obv{ordinal_elem_implies_subseteq}.
    Follows by \auto{pno_lt_initial,subseteq}.
\end{proof}

\begin{theorem}\label{pno_lt_initial_trichotomy}
    Let $\alpha$ be an ordinal.
    Then $\forall p,q:\Set\to\Prop.
        \Apply{\pnoLtInitialOp}{\alpha,p,q}\lor
        \Apply{\pnoEqOp}{\alpha,p,q}\lor
        \Apply{\pnoLtInitialOp}{\alpha,q,p}$.
\end{theorem}
\begin{proof}
    Fix $p,q$.
    \begin{byCase}
        \caseOf{$\Apply{\pnoEqOp}{\alpha,p,q}$.}
            Follows by \obv{pno_eq,pno_lt_initial}.
        \caseOf{$\lnot\Apply{\pnoEqOp}{\alpha,p,q}$.}
            Let $D=\sep{\alpha}{\fun{\delta:\Set}{
                (\Apply{p}{\delta}\land\lnot\Apply{q}{\delta})\lor
                (\lnot\Apply{p}{\delta}\land\Apply{q}{\delta})}}$.
            $D$ is inhabited
                by \auto{pno_eq,logical_separation_iff}.
            Take $\delta$ such that $\delta$ is an \in-minimal element of $D$
                by \obv{regularity}.
            $\delta\in D$ and $\delta\notmeets D$.
            Hence $\delta\in\alpha$
                by \obv{logical_separation_iff}.
            Hence
                $((\Apply{p}{\delta}\land\lnot\Apply{q}{\delta})\lor
                 (\lnot\Apply{p}{\delta}\land\Apply{q}{\delta}))$
                by \obv{logical_separation_iff}.
            Show that
                $\Apply{\pnoEqOp}{\delta,p,q}$.
            \begin{subproof}
                Show that $\forall \gamma:\Set.\ \gamma\in\delta\implies
                    (\Apply{p}{\gamma}\iff\Apply{q}{\gamma})$.
                \begin{subproof}
                    Fix $\gamma\in\delta$.
                    Hence $\gamma\in\alpha$
                        by \obv{ordinal_transitivity}.
                    Suppose not.
                    Hence $\gamma\in D$
                        by \auto{logical_separation_iff}.
                    Contradiction by \auto{disjoint}.
                \end{subproof}
                Follows by \obv{pno_eq}.
            \end{subproof}
            \begin{byCase}
                \caseOf{$\Apply{p}{\delta}$.}
                    We have $\lnot\Apply{q}{\delta}$.
                    We have
                        $\Apply{\pnoEqOp}{\delta,q,p}$
                        by \obv{pno_eq_symmetric}.
                    We have
                        $\Apply{\pnoLtInitialOp}{\alpha,q,p}$
                        by \obv{pno_lt_initial_intro}.
                \caseOf{$\lnot\Apply{p}{\delta}$.}
                    We have $\Apply{q}{\delta}$.
                    We have
                        $\Apply{\pnoLtInitialOp}{\alpha,p,q}$
                        by \obv{pno_lt_initial_intro}.
            \end{byCase}
    \end{byCase}
\end{proof}
